\section{Heuristic interpretation, literature and limitations}\label{sec:heuristic}
\subsection{Proposed mechanism, explicitly heuristic}
Proposition~\ref{prop:row} interprets each matrix row as a
\emph{normalized shifted sextic Gauss-sum first moment}
with shift \(u=\rho\kappa\). It is plausible that
averaging these moments across row primes produces a
crossover when the conductor scale of the shift changes
relative to the prime-modulus cutoff. This is a proposed
mechanism, \emph{not} derived from the Kubota functional
equation or from known theta-coefficient formulas.

The large-\(t\) fifth-power deficit is compatible with
the vanishing condition for some sextic theta
coefficients at fifth prime powers discussed by
Br\"oker--Hoffstein \cite{broeker2013};
compatibility does not determine the \emph{signed}
matrix energy. Similarly, the exploratory exponent-two
row concentration resembles arithmetic local-factor
selection, without constituting a proof.

The scale correction in
\texttt{reciprocity\_spectra/RS\_HP\_SCALE\_ERRATUM.md}
is essential: a Patterson-type sextic prime moment
at size \(X^{2/3}/\log X\) is an asymptotic for
\emph{normalized} sums \(g_\sigma(u)/\sqrt{\Nm\sigma}\).
The analogous raw Gauss-sum scale would be
\(X^{7/6}/\log X\). The former raw-scale
hypothesis in RS-S026/S027 was an error and
its dependent rates must not be invoked as
results for the corrected model.

Even a uniform \(o(X^{2/3}/\log X)\) error in
\(H_P\) exceeds the typical row size
\(\sqrt D\), so it provides no direct control
of the signed excess in
Proposition~\ref{prop:gram}.
A viable proof would require substantially
finer, row-averaged \emph{signed two-point}
estimates in the moving-shift range.
The programme's mathematical proof work on
that question is parked rather than
misrepresented as completed.

\subsection{Relation to existing work}
Residue-symbol matrices and their
characterizations were studied by
Dummit--Dummit--Kisilevsky \cite{ddk2016}.
Patterson, Heath-Brown--Patterson,
Kazhdan--Patterson, and
Chinta--Friedberg--Hoffstein supply part
of the metaplectic/character-sum framework
\cite{patterson1987,hbp1979,kazhdan1984,chinta2003}.
Br\"oker--Hoffstein analyse quartic and
sextic metaplectic coefficients
\cite{broeker2013}. For cubic Gauss sums,
Dunn--Radziwi\l{}\l{} investigate
Patterson's prime-argument bias
\cite{dunn2024}, while David--Dunn--Hamieh--Lin
treat the quartic setting \cite{ddhl2023}.
Blomer--Goldmakher--Louvel provide
large-sieve context \cite{bgl2014};
de Faveri--Dunn--Hoffstein study
nonorthogonality phenomena
\cite{dfdh2026}; Xu--Sang provide recent
norm-level sextic context \cite{xusang2026}.
The analyses here neither claim
independence from these papers nor
infer novelty from a literature search.

\subsection{Limitations and failed antecedents}
The main empirical limitation is that all
arithmetic measurements end at
\(X\le4\cdot10^6\); no data here establish
the function \(F\), uniform limiting
convergence or a fixed-class threshold.
Some plotted points share class lists and
engines; error bars based solely on
independent sampling would be misleading.

The universal norm-only predecessor
CAND-001 is \emph{refuted} by sixth-power
invariance at \(u=64\), and the attempted
repair PROV-RS14 was \emph{refuted} by its
registered RS-S015 class tests.
The per-class fifth-power clause of
RS-S017 Q3 \emph{failed}. The attractive
finite mean deficit is not a justification
for erasing any of these failures.
The H\(_P\) raw-versus-normalized erratum,
the incorrect historical logical-corollary
wording and the absent signed prime-correlation
estimate are further substantive limitations.

\subsection{Availability and responsible-use disclosure}
The code, numerical records, frozen predictions,
session ledger and this manuscript source are in
\url{https://github.com/jfairfaxball-348/PRIMES}.
The primary data are the committed machine-readable
files under \texttt{reciprocity\_spectra/results/};
the plotting program is
\texttt{paper/figures/make\_all.py}. The finite Lean
package is \texttt{formalisation/} on the pinned
toolchain; independently verified selected theorems
are listed in its README. Each plot is reproducible
from the committed JSON, and conceptual diagrams
are labelled as such. No Palomar ID is claimed.

\textbf{AI assistance.}
OpenAI's ChatGPT was used substantially in the
Lean formalisation and assisted with manuscript
drafting. The author is responsible for the
mathematical claims, numerical evidence,
verification, presentation and submission.
The exact underlying model versions were not
consistently recorded. Machine checking of the
finite Lean results does not independently
validate the conjectures or this written text.


